%
% Halaman Abstrak
%
\setstretch{1}
\vspace*{0.2cm}
{
    \singlespacing
    \setlength{\parindent}{0pt}

    \begin{tabular}{@{}l l p{10cm}}
        Nama&: & \penulis \\
        Program Studi&: & \program \\
        Judul&: & \judul \\
        Pembimbing&: & \pembimbing \\
    \end{tabular}

    \bigskip
    \bigskip

\textit{Vision-language model} (VLM) dan \textit{vision-language-action} (VLA) berpotensi digunakan untuk perencanaan semantik serta pengusulan aksi dalam sistem navigasi UAV. Namun, perbedaan prapemrosesan masukan, templat \textit{prompt}, format keluaran, dan ruang aksi menyulitkan penggantian model serta perbandingan yang adil dalam sistem navigasi \textit{closed-loop}. Penelitian ini mengusulkan DroneVL, yaitu adaptor VLM dan VLA agnostik-model untuk navigasi UAV semantik di CoSyS-AirSim yang terintegrasi dengan ROS~2. Eksperimen VLM membandingkan GPT-5.6, Qwen3.6 35B-A3B, dan Gemini Robotics ER 2 dengan varian \textit{Streaming Preview}. Eksperimen portabilitas VLA membandingkan CognitiveDrone, OpenVLA-7B yang diadaptasi untuk UAV, dan OpenPI $\pi_{0.5}$ yang diadaptasi untuk UAV; dua model manipulasi terakhir dilatih menggunakan trajektori CoSyS-AirSim. DroneVL menormalkan citra RGB, data kedalaman, keadaan UAV, dan instruksi misi ke dalam kontrak pengamatan bersama. Setiap respons atau keluaran aksi diurai, dinormalisasi, dan divalidasi menjadi aksi navigasi tingkat tinggi sebelum dijalankan oleh pengendali penerbangan tingkat rendah yang tetap. Eksperimen utama membandingkan model dasar OpenVLA-7B, OpenVLA-7B dengan LoRA yang diadaptasi untuk UAV, dan kondisi LoRA yang dilanjutkan GRPO pada \textit{rollout} simulator. Penelitian ini juga merancang DroneVL Arena, yaitu kerangka \textit{benchmark} dengan episode misi berversi, anotasi semantik, batasan keselamatan, pencatatan lintasan, dan evaluator. Evaluasi mencakup penambatan visual, penalaran spasial, pencarian objek, navigasi semantik, pelaksanaan instruksi multilangkah, keselamatan, efisiensi lintasan, aksi tidak valid, dan latensi. Luaran yang diharapkan adalah lapisan integrasi dan \textit{benchmark} yang dapat direproduksi untuk menganalisis portabilitas dan adaptasi \textit{closed-loop} VLM/VLA pada UAV simulasi.

    \bigskip

    \textbf{Kata kunci:} \textit{vision-language model}, \textit{vision-language-action}, UAV, navigasi semantik, CoSyS-AirSim, ROS~2, LoRA, GRPO, \textit{benchmark}.
}

\newpage
